% Part: sets-functions-relations
% Chapter: arithmetization
% Section: reflections

\documentclass[../../../include/open-logic-section]{subfiles}

\begin{document}

\olfileid{sfr}{arith}{ref}
\olsection{Enkele filosofische beschouwingen}
Tot zover de technische details. Maar wat hebben ze opgeleverd?
In elk geval leverden ze {fraaie} zuivere wiskunde op. Daarnaast waren
er diepgaande begripsmatige resultaten. Het inzicht dat de
volledigheidseigenschap het wezenlijke verschil tussen de reële en de
rationale getallen uitdrukt, was diepgaand. Bovendien geeft de
expliciete constructie van de reële getallen als Dedekindsneden de
analyse een stevige grondslag. We weten dat een \emph{volledig
geordend lichaam} kan bestaan, want de sneden vormen zo'n lichaam.
Toch zijn bij deze resultaten enkele kanttekeningen op hun plaats.
Ten eerste is het niet duidelijk dat reële getallen als sneden
opvatten \emph{strenger} is dan ze opvatten via hun vertrouwde,
mogelijk oneindige decimaalontwikkelingen. Die laatste ``constructie''
van de reële getallen lijkt enigszins op de constructie met
Cauchyrijen, maar was in wezen al vanaf het begin van de zeventiende
eeuw bij wiskundigen bekend (zie \olref[cauchy]{sec}). De werkelijke
toename in strengheid kwam door het inzicht dat de reële getallen de
volledigheidseigenschap hebben. Dat we reële getallen als bepaalde
verzamelingen kunnen construeren, is op zichzelf wellicht niet zo
interessant.
Nog minder duidelijk is dat de veel eenvoudigere aritmetisering van de
gehele of rationale getallen de wiskundige behandeling ervan strenger maakt.
Hierbij is een eenvoudige vaststelling van belang. Nadat we de gehele
getallen als equivalentieklassen van geordende paren natuurlijke
getallen hebben \emph{geconstrueerd}, vervolgens de rationale getallen
als equivalentieklassen van geordende paren gehele getallen en ten
slotte de reële getallen als verzamelingen rationale getallen, gaan we
die constructies onmiddellijk weer \emph{vergeten}. Niemand zal tijdens
een wiskundig bewijs op deze constructies willen \emph{teruggrijpen},
behalve natuurlijk in een bewijs dat ze zich gedragen zoals bedoeld.
Het is veel eenvoudiger om rechtstreeks over een reëel getal te
spreken dan over een verzameling van verzamelingen van verzamelingen
van verzamelingen van verzamelingen van verzamelingen van verzamelingen
natuurlijke getallen.
%(For reals were constructed as sets of rationals; and these were constructed as equivalence classes of ordered pairs of integers, i.e., sets of sets of sets of integers; etc.).
%2/3 = [2,3]
%	i.e. {<2,3>, ... }
%	{ {{2},{2,3}}, ... }
%
%reals are 
%	sets of rationals 
%
%rationals are 
%	sets of sets of sets of integers
%
%integers are
%	sets of sets of sets of naturals
Het twijfelachtigst is dat deze definities ons vertellen wat de gehele,
rationale of reële getallen \emph{zijn}, \emph{metafysisch gezien}.
Het is dus twijfelachtig of de reële getallen bijvoorbeeld
\emph{identiek zijn aan} bepaalde verzamelingen van verzamelingen van
verzamelingen\ldots.
Het voornaamste bezwaar is dat de constructie op veel verschillende
manieren had kunnen worden uitgevoerd. Voor de reële getallen bestaan
enkele werkelijk interessante alternatieve constructies (zie
\olref[cauchy]{sec}). Maar er is ook een volkomen triviale manier om
andere constructies te krijgen. Zoals in
\olref[sfr][rel][ref]{sec} hadden we geordende paren iets anders kunnen
definiëren. Met dat alternatieve begrip van een geordend paar zouden
onze constructies even goed hebben gewerkt, maar andere objecten hebben
opgeleverd. Er bestaan dus verscheidene concurrerende constructies
binnen de verzamelingenleer van de gehele, rationale en reële
getallen. Het zou willekeurig en gênant zijn om te beweren dat de
gehele getallen bijvoorbeeld \emph{deze} verzamelingen zijn en niet
\emph{die}. (Net als in \olref[sfr][rel][ref]{sec} is dit een geval van
een argument dat door \citealt{Benacerraf1965} bekend is geworden.)
Er is nog een punt: onze constructies hebben iets uitgesproken
\emph{vreemds}. We begonnen met de natuurlijke getallen. Daarna
construeren we de gehele getallen, waaronder ``de $0$ van de gehele
getallen'', namelijk $ \equivrep{0,0}{\Intequiv}$. Maar $0 \neq
\equivrep{0,0}{\Intequiv}$. Volgens onze constructies is zelfs
\emph{geen enkel} natuurlijk getal een geheel getal. Dat lijkt zeer
tegenintuïtief. In \olref[sfr][set][imp]{sec} beweerden we immers
zonder veel toelichting dat $\Nat \subseteq \Rat$. Als de constructies
ons precies vertellen \emph{wat} de getallen zijn, was die bewering
evident onwaar.
Wat is dan de slotsom? Binnen een na\"ieve verzamelingenleer, waarin we
de natuurlijke getallen als gegeven aannemen, kunnen we gehele,
rationale en reële getallen als bepaalde verzamelingen
\emph{behandelen}. In die zin kunnen we de theorieën van deze objecten
in een verzamelingenleer \emph{inbedden}. De filosofische betekenis van
zo'n inbedding is echter niet zonder meer duidelijk.
Dit is natuurlijk niet het laatste woord!{} Het punt is alleen dit. Om
aan te tonen dat de aritmetisering van de reële getallen werkelijk van
groot filosofisch belang \emph{is}, is nog een aanvullend
\emph{filosofisch} argument nodig.
%Additionally: for the entire duration of this chapter, we have treated
%the natural numbers as simply \emph{given} to us. But what can we do
%with them? That is the subject matter for the next chapter.
\end{document}
